\documentclass{beamer}
\usetheme{Madrid}
\usepackage[utf8]{inputenc}
\title{Clase 5 — Introducción a Redes Neuronales}
\author{Curso: Laboratorio IA}
\date{Semana 5}

\begin{document}
\frame{\titlepage}

\begin{frame}{Objetivo}
  Introducir redes neuronales con un ejemplo práctico en PyTorch: predecir la estatura en niños a partir de su edad.
\end{frame}

\begin{frame}{Agenda}
  \begin{itemize}
    \item Motivación y problema
    \item Dataset sintético (CSV)
    \item Modelo simple en PyTorch
    \item Entrenamiento y evaluación
    \item Demo y conclusiones
  \end{itemize}
\end{frame}

\begin{frame}{Motivación}
  \begin{itemize}
    \item Las redes neuronales aprenden relaciones no lineales entre variables.
    \item Empezamos con un problema sencillo: edad $\rightarrow$ estatura.
    \item Usamos datos sintéticos para ilustrar flujo de trabajo (limpio y reproducible).
  \end{itemize}
\end{frame}

\begin{frame}{Dataset sintético}
  \begin{itemize}
    \item Archivo: data/estatura_ninos.csv (columnas: `age`, `height`).
    \item Generado con ruido para simular variabilidad real.
    \item Usamos $\sim$500 muestras entre 0 y 18 años.
  \end{itemize}
\end{frame}

\begin{frame}{Modelo propuesto}
  \begin{itemize}
    \item Red neuronal pequeña: capas lineales + ReLU.
    \item Entrada: edad (normalizada), salida: estatura (cm).
    \item Pérdida: MSE. Optimizador: Adam.
  \end{itemize}
\end{frame}

\begin{frame}{Entrenamiento}
  \begin{itemize}
    \item División entrenamiento / prueba.
    \item Monitorizar pérdida y error medio absoluto (MAE).
    \item Guardar modelo y visualizar predicciones vs. reales.
  \end{itemize}
\end{frame}

\begin{frame}{Código de ejemplo}
  \begin{itemize}
    \item Ver `scripts/train_estatura.py` para el código PyTorch completo.
    \item Ver `data/generate_height_data.py` para reproducir el CSV.
  \end{itemize}
\end{frame}

\begin{frame}{Tareas y próximas sesiones}
  \begin{itemize}
    \item Ejecutar el entrenamiento en clase y discutir resultados.
    \item Extender a múltiples características (peso, sexo, etc.).
    \item Introducir regularización y validación cruzada.
  \end{itemize}
\end{frame}

\end{document}
